
% =============================================================================
% 1. INTRODUCTION
% =============================================================================
\section{Introduction}
\label{sec:intro}
Early identification of foliar plant diseases is critical for timely agricultural intervention, reducing crop losses and preventing unnecessary pesticide application~\cite{richard2022control}. Because foliar lesions can spread rapidly across agricultural canopies, manual scouting by agronomists is labor-intensive, subjective, and difficult to sustain across large fields~\cite{shoaib2023advanced,radocaj2025evolution}. Automated vision systems offer a scalable alternative, but their heavy reliance on large labeled datasets remains a primary limitation~\cite{shoaib2023advanced}.

Deep learning architectures, including convolutional neural networks (CNNs) and vision transformers (ViTs), achieve strong performance on standardized crop pathology benchmarks~\cite{gupta2024fruit,dinesh2025plant}. However, conventional supervised models require large labeled datasets~\cite{shoaib2023advanced}. In agricultural practice, new pathogen variants, regional strains, and early-stage lesions provide only scarce and imbalanced data. As a result, standard training pipelines adapt poorly to rapidly evolving field conditions.

Few-shot learning (FSL) reduces this dependence by adapting to novel classes from a few labeled examples per pathology~\cite{sun2024few}. Metric-based methods compare query leaves with support examples in a learned metric space~\cite{zhao2024mafde,snell2017prototypical}. Standard distances such as Euclidean or cosine treat feature coordinates uniformly and ignore class-conditional covariance~\cite{bateni2020improved}; background texture and variable illumination add further error in field images~\cite{yang2024laboratory}. Parameter-efficient fine-tuning (PEFT) and optimization-based meta-learning can learn sharper boundaries, but they depend on iterative backpropagation and optimizer state~\cite{hu2022lora,finn2017maml}, which exceed the memory and power budgets of agricultural edge hardware~\cite{radocaj2025evolution}.

Edge diagnosis methods have two practical gaps. On-device cost is rarely measured: the computing requirements of field models remain under-researched and efficiency metrics are reported inconsistently~\cite{radocaj2025evolution}. In addition, standard few-shot classifiers return uncalibrated softmax scores that do not convey predictive uncertainty, which can produce overconfident errors on ambiguous inputs~\cite{guo2017calibration}. Without uncertainty-aware routing, automated scouting can lead to incorrect treatment or missed infections.

We therefore developed BioDINO-Edge, a gradient-free, two-tier cascaded few-shot framework for edge execution. It couples a fast Tier-1 visual encoder (8-bit integer (INT8) quantized DINOv2-Small~\cite{oquab2023dinov2}) with a deeper Tier-2 biological encoder (INT8-quantized BioCLIP~\cite{stevens2024bioclip}) in a joint 896-dimensional space. A dual-metric Shannon entropy and margin gate evaluates Tier-1 confidence, so clear leaves exit early and uncertain queries move to the deeper Tier-2 model. Adaptation to novel pathologies is computed analytically with regularized Fisher Linear Discriminant Analysis (Fisher-LDA)~\cite{fisher1936use,rao1948utilization} and Ledoit--Wolf covariance shrinkage~\cite{ledoit2004well}; one-time adaptation takes 0.53\,s on a mobile central processing unit (CPU) without backpropagation. The complete system runs as a standalone Android application package (APK) and executes fully offline with low memory overhead.

The primary contributions of this work are:
\begin{enumerate}
    \item \textbf{Closed-form, gradient-free edge adaptation:} We formulate few-shot foliar adaptation as a closed-form algebraic solution using analytical Ledoit--Wolf covariance shrinkage and regularized Fisher-LDA projections. This enables novel foliar pathologies to be registered on resource-constrained mobile CPUs in 0.53\,s without iterative backpropagation, autograd graphs, or optimizer tuning.
    \item \textbf{Uncertainty-gated cascaded inference:} We design a dual-metric confidence gate (normalized Shannon entropy and top-2 margin) that routes clear foliage to a low-latency Tier-1 fast exit while reserving multimodal consensus for ambiguous queries. By routing 14.02\% of queries through the 59.78\,ms fast path, the cascade reduces average inference latency from 164.52\,ms to 149.84\,ms while preserving 97.41\% accuracy. Transductive recalibration yields an additional +2.33 percentage points in accuracy without gradient computation.
    \item \textbf{Physical edge benchmarking:} We test dual-foundation inference on an ARM64 mobile CPU without GPU acceleration, evaluating end-to-end latency, memory footprint ($<$500\,MB PSS), and sustained continuous throughput (settling at 4.26 frames per second [FPS] as the Android thermal status rises from 2 to 3 during a 600\,s soak).
    \item \textbf{Field-domain evaluation with target-domain support:} On PlantDoc, combining the botanical text prior with regularized visual features over the same frozen encoders reaches 84.94\% $\pm$ 11.15\% in the 5-way 5-shot protocol and 52.38\% across all 27 classes, with support drawn from the PlantDoc training split and no backbone fine-tuning.
\end{enumerate}

The remainder of this paper is organized as follows. Section~2 reviews related work in agricultural few-shot learning, transductive inference, and edge cascades. Section~3 formulates the sequential edge few-shot task, derives the operational constraints, and analyzes small-sample covariance degeneracy. Section~4 details the proposed BioDINO-Edge methodology, including the dual-foundation representation, closed-form regularized Fisher-LDA adaptation, and the two-tier uncertainty cascade. Section~5 describes the experimental setup, benchmark datasets, and hardware profiling testbed. Section~6 reports benchmark comparisons, statistical significance tests, component ablations, computational latency, and robustness evaluations. Section~7 discusses operational deployment considerations and practical limitations, and Section~8 concludes the paper with future directions.
